\section{外层临界极点与显式弱型界}\label{sec:outer-critical-app}
\subsection{正权的临界系数与有限性}
取 $\nu=dx$、$d\mu=w\,dx$，其中 $w>0$ a.e. 且 $w\in L^1_{\rm loc}$。
记 Lebesgue 平均的加权弱型范数为
\[W^{\rm out}_{K,w}=\|M_K\|_{L^1(w)\to L^{1,\infty}(w)}.
\tag{OW.1}\]
定理 \ref{thm:two-measure-weak} 及推论 \ref{cor:two-measure-pole} 给出
\[\lim_{\alpha\downarrow0}\alpha C^{\rm out}_{K,w}(\alpha)=W^{\rm out}_{K,w},\qquad
\lim_{\lambda\downarrow\lambda_c}(\lambda-\lambda_c)
C^{\boldsymbol a,\rm out}_{K,w}(\lambda)
=\rho_c\lambda_c W^{\rm out}_{K,w},\tag{OW.2}\]
允许扩展值 $+\infty$。若该弱型范数有限，超临界常数均有限；临界及以下发散。

\begin{proposition}\label{prop:outer-finite-weak}
若 $w$ 满足 (H.1)，则
\[W^{\rm out}_{K,w}\le B_K\frac ba\frac n\tau3^n
\le\frac ba\frac n\tau9^n<\infty.\tag{OW.3}\]
这里 $B_K$ 可用 (G.1) 的测度一致弱型上界；方体可取 $B_K=2^n$。
此外，对本节任意正权，
\[W^{\rm out}_{K,w}\ge W_{K,dx}.\tag{OW.4}\]
\end{proposition}
\begin{proof}
令 $\gamma=n-\tau$、$Q=x+rK$。若 $r\le\rho(x)/2$，两侧比较给出
$\operatorname*{ess\,sup}_Qw/\operatorname*{ess\,inf}_Qw\le(b/a)3^\gamma$。若 $r>\rho(x)/2$，则
\[\frac{\mu(Q)}{|Q|\operatorname*{ess\,inf}_Qw}
\le\frac ba\frac n\tau\left(\frac{\rho(x)+r}{r}\right)^n
\le\frac ba\frac n\tau3^n.\]
因此令 $A=(b/a)(n/\tau)3^n$，在每个球上均有
$\int_Q f/|Q|\le A\int_Q f\,d\mu/\mu(Q)$，故
$M_Kf\le A M_{K,\mu}f$。应用测度一致弱型界即得 (OW.3)。

为证 (OW.4)，取 $w$ 的正有限 Lebesgue 点 $x_0$。
固定有界紧支撑 $f\ge0$、$\ell>0$ 与紧集 $L\subset\{M_Kf>\ell\}$。
令 $f_\varepsilon(x)=f((x-x_0)/\varepsilon)$。球族的伸缩不变性给出
$x_0+\varepsilon L\subset\{M_Kf_\varepsilon>\ell\}$，而 Lebesgue 微分定理给
\[\varepsilon^{-n}\int f_\varepsilon w\to w(x_0)\int f,\qquad
\varepsilon^{-n}\mu(x_0+\varepsilon L)\to w(x_0)|L|.\]
输入误差由 $\|f\|_\infty$ 乘以固定有界区域上的密度 $L^1$ 误差控制。
取 $f,\ell,L$ 的上确界，即得 (OW.4)。
\end{proof}

对全部正权，超临界有限性不能无条件成立。例如一维取 $w(x)=e^x$、
$E_L=(0,L)$。若 $0<\alpha<1/2$，对
$L<x<L/(2\alpha)$ 取中心区间 $[0,2x]$，其平均 $L/(2x)>\alpha$。因此
\[C^{\rm out}_{c,e^x}(\alpha)\ge
\frac{e^{L/(2\alpha)}-e^L}{e^L-1}\longrightarrow\infty.
\tag{OW.5}\]
当 $m=2$、$a_1=a_2=1$ 时，$1/4<\lambda<9/16$ 对应
$0<\Phi_2(\lambda)<1/2$，故这些超临界点的外层常数仍为无穷。

\subsection{一维幂权的构造下界与比较上界}
对 $w_\tau(x)=|x|^{\tau-1}$、$0<\tau<1$，记
$W_\tau=W^{\rm out}_{c,w_\tau}$，并令 $u_\tau\in(0,1)$ 是唯一解
\[u_\tau^{1-\tau}+(1-\tau)u_\tau=\tau.\tag{OW.6}\]
\begin{proposition}\label{prop:atomic-power-bound}
有严格可复核的界
\[\max\left\{\frac{11+\sqrt{61}}{12},
\frac3{4\tau}\left(\frac53\right)^\tau\right\}
\le W_\tau\le2u_\tau^{\tau-1}.\tag{OW.7}\]
特别地，
\[\frac{\sqrt{15}}2\le W_{1/2}\le2(1+\sqrt2),\qquad
W_{1/4}\ge3(5/3)^{1/4}>\frac{17}5.\tag{OW.8}\]
\end{proposition}
\begin{proof}
第一个下界由 (OW.4) 及 Melas 的一维 Lebesgue 弱型常数
\cite{Melas2003} 得到。
保持 Lebesgue 平均，先以输入测度
$\sigma=\frac43(\delta_{-1}+\delta_1)$ 计算
$M_c\sigma(x)=\sup_{r>0}\sigma([x-r,x+r])/(2r)$。
单原子与双原子的见证区间给出严格光晕
\[\{M_c\sigma>1\}=(-5/3,-1/3)\cup(-1/3,1/3)\cup(1/3,5/3).\]
这些区间的加权总质量为 $2(5/3)^\tau/\tau$，输入成本为 $8/3$。
以 $\frac{4/3}{2\varepsilon}\1_{[x_j-\varepsilon,x_j+\varepsilon]}$ 替换各原子。
权在 $x_j=\pm1$ 连续，故成本趋于 $8/3$；每个严格原子见证都可微增半径，
使平滑输入仍具有该严格见证。Fatou 引理保存光晕质量下界，
因此 $W_\tau\ge\frac3{4\tau}(5/3)^\tau$。

为证上界，将任意区间反射、伸缩，使绝对值最大的端点为 $1$，另一端为
$v\in[-1,1]$。其平均权与本质下确界之比，在 $v\ge0$ 时至多为 $1/\tau$，
在 $v=-u\le0$ 时等于
\[g(u)=\frac{1+u^\tau}{\tau(1+u)}.
\]
导数的符号等于 $\tau-u^{1-\tau}-(1-\tau)u$ 的符号，后者严格递减，
故最大值在 (OW.6) 的唯一解取得，且为 $A_\tau=u_\tau^{\tau-1}$。
逐区间比较给 $M_cf\le A_\tau M_{c,w_\tau}f$；后者的测度一致弱型界为 $2$，
从而 $W_\tau\le2A_\tau$。
当 $\tau=1/2$，有 $u_\tau=(\sqrt2-1)^2$。最后，
$[3(5/3)^{1/4}]^4-(17/5)^4=854/625>0$，给出 (OW.8)。
\end{proof}

例如对 $m=2$、等指数、$w=|x|^{-3/4}$，有
\[\lim_{\lambda\downarrow1/4}(\lambda-1/4)
C^{(2),\rm out}_{c,w}(\lambda)
=\tfrac12W_{1/4}\ge\tfrac32(5/3)^{1/4}>\frac{17}{10}.\tag{OW.9}\]
同一权的内禀模型由 $W_\mu\le2$ 得临界系数至多为 $1$。
这些界分别针对各自的平均测度；(OW.7) 尚未确定 $W_\tau$ 的精确值。
